\section{\MakeUppercase{1. The shortcut}}

\lead{On BRENDA-200 (200 proteins\,$\times$\,200 molecules, protein-stratified
split) four published models reach a \hlw{pooled AUC of 0.63--0.95}. Ask them
the question the task actually poses, \emph{for this molecule, put its true
targets on top}, and all four collapse to chance.}

\panel{\panelhead{How we score}
{\footnotesize \textbf{Pooled AUC}: one number over all 40,000 pairs at once.
\textbf{Matrix MRR}: within one molecule's row of the 200\,$\times$\,200 score
matrix, the mean of 1/rank of its true target; 1.0 = always first,
\textbf{0.029} = chance. \textbf{Hit@10}: true target in the top ten.
\textbf{Cheat sheet}: molecule-blind baseline scoring each protein by its
training positive rate; \emph{overlap} = Jaccard of the top-ten sets,
\emph{concentration} = Gini of the scores.}}

\vfill

{\small Why? Because the cheat sheet does just as well: it picks the same
favourites, up to two thirds of their top ten, at the same concentration
($\approx$ 0.995). Pooled AUC certifies the protein, not the pair. We retire it
as a success gate and rank inside each molecule instead.}

\vfill

\includegraphics[width=\linewidth]{fig_dissociation.pdf}
\figcap{\textbf{The two scores disagree.} Horizontal: what the field reports.
Vertical: ranking AUC within one molecule, chance 0.5.}

\vfill

{\footnotesize
\begin{tabular}{@{}l r r r r@{}}
	\toprule
	Model & pooled & one-mol. & cheat & conc. \\
	      & AUC    & AUC      & overlap & \\
	\midrule
	DrugBAN         & 0.954 & \textcolor{rbwarm}{0.433} & 0.54 & 0.995 \\
	MolTrans        & 0.937 & \textcolor{rbwarm}{0.470} & 0.05 & 0.987 \\
	GraphDTA        & 0.869 & \textcolor{rbwarm}{0.533} & 0.67 & 0.995 \\
	GEMS            & 0.633 & \textcolor{rbwarm}{0.514} & 0.67 & 0.995 \\
	\midrule
	\textit{cheat sheet} & n/a & n/a & 1.00 & 0.995 \\
	\midrule
	\textbf{RankBind} & 0.634 & \hlb{0.891} & \textbf{0.035} & \textbf{0.787} \\
	\bottomrule
\end{tabular}}
\figcap{Chance is 0.50 for one-molecule AUC; three of four baselines rank
\emph{below} it. \textit{n} = 30 held-out positives.}

\vfill

\includegraphics[width=\linewidth]{fig_respmaps.pdf}
\figcap{\textbf{What the shortcut looks like.} 200 molecules (rows)\,$\times$\,200
proteins (columns), darker = higher. A vertical stripe is a protein that wins
almost every molecule. GraphDTA's map is the cheat sheet's; RankBind's is
de-striped.}
